% (c) Nikita Lisitsa, lisyarus@gmail.com, 2026

\documentclass[10pt]{beamer}

\usepackage[T2A]{fontenc}
\usepackage[russian]{babel}
\usepackage{minted}

\usepackage[most]{tcolorbox}
\tcbuselibrary{minted}

\usepackage{graphicx}
\graphicspath{ {./images/} }

\usepackage{adjustbox}

\usepackage{color}
\usepackage{soul}

\usepackage{hyperref}

\usepackage{tikz}
\usetikzlibrary{decorations}
\usetikzlibrary{decorations.pathreplacing}
\usepackage{xifthen}

\usepackage{alltt}

\usetheme{metropolis}
\setminted{fontsize=\footnotesize}

\definecolor{red}{rgb}{1,0,0}
\definecolor{green}{rgb}{0,0.5,0}
\definecolor{blue}{rgb}{0,0,1}
\definecolor{magenta}{rgb}{0.75,0,0.75}
\definecolor{codebg}{RGB}{29,35,49}

\makeatletter
\newcommand{\slideimage}[1]{
  \begin{figure}
    \begin{adjustbox}{width=\textwidth, totalheight=\textheight-2\baselineskip-2\baselineskip,keepaspectratio}
      \includegraphics{#1}
    \end{adjustbox}
  \end{figure}
}
\makeatother

\usemintedstyle{tango}

\title{Компьютерная графика}
\subtitle{Лекция 5: Bind groups, uniform буферы, текстуры, сэмплеры}
\date{2025}

\setbeamertemplate{footline}[frame number]

\begin{document}

\frame{\titlepage}

\begin{frame}[fragile]
\frametitle{Как мы умеем передавать данные на GPU?}
\begin{itemize}
\item Immediates -- ограничен размер
\pause
\item Атрибуты вершин, индексы -- ограничены способы использования (только через растеризацию, строгий формат, нет доступа к соседним вершинам)
\pause
\item Хочется уметь передавать больше произвольных данных
\pause
\item \(\Longrightarrow\) \textit{bind группы}
\end{itemize}
\end{frame}

\begin{frame}[fragile]
\frametitle{Bind groups}
\begin{itemize}
\item Bind group -- контейнер объектов (\textit{bindings}), которые можно передать на GPU в шейдер
\pause
\item Могут содержать:
\begin{itemize}
\item Буферы: \textit{uniform}, \textit{storage}, \textit{read-only storage}
\pause
\item Текстуры (read-only изображения в памяти GPU)
\pause
\item Сэмплеры (настройки чтения изображений)
\pause
\item Storage текстуры (read-write изображения)
\end{itemize}
\end{itemize}
\end{frame}

\begin{frame}[fragile]
\frametitle{Bind groups}
\begin{itemize}
\item Объекты в bind группе описываются по числовым индексам
\pause
\item Пример bind-группы:
\end{itemize}
\usemintedstyle{lightbulb}
\begin{tcblisting}{listing engine=minted, minted language=cpp, minted options={fontsize=\scriptsize}, listing only, colback=codebg, colframe=codebg, rounded corners}
binding 0: uniform-буфер с параметрами камеры
binding 1: environment-текстура для отражений
binding 2: sampler для чтения текстуры отражений
binding 3: буфер с описанием источников света
\end{tcblisting}
\end{frame}

\begin{frame}[fragile]
\frametitle{Bind group layout}
\begin{itemize}
\item Для создания bind group нужно сначала описать \textit{bind group layout}
\pause
\item Layout описывает, объекты какого типа (буфер/текстура/etc) будут храниться в bind группе (без ссылки на конкретные объекты), и в каких шейдерах они будут использоваться (вершинный/фрагментный)
\pause
\item Пример bind group layout'а:
\end{itemize}
\usemintedstyle{lightbulb}
\begin{tcblisting}{listing engine=minted, minted language=cpp, minted options={fontsize=\scriptsize}, listing only, colback=codebg, colframe=codebg, rounded corners}
binding 0: uniform-буфер
binding 1: 2D текстура
binding 2: sampler
binding 3: storage буфер
\end{tcblisting}
\end{frame}

\begin{frame}[fragile]
\frametitle{Pipeline layout}
\begin{itemize}
\item При создании render pipeline нужно указать, какие ресурсы нужны этому pipeline для рисования
\pause
\item Это делается через создание \textit{pipeline layout}: он описывает набор bind group layout'ов, которые ему нужны
\pause
\item Сами bind группы внутри pipeline layout тоже описываются индексами
\pause
\item Пример pipeline layout'а:
\end{itemize}
\usemintedstyle{lightbulb}
\begin{tcblisting}{listing engine=minted, minted language=cpp, minted options={fontsize=\scriptsize}, listing only, colback=codebg, colframe=codebg, rounded corners}
group 0: настройки камеры и освещения
group 1: параметры материала объекта
group 2: вспомогательные текстуры (шум, дизеринг)
\end{tcblisting}
\end{frame}

\begin{frame}[fragile]
\frametitle{Bind groups + pipelines вместе}
\begin{itemize}
\item \textit{Bind group / pipeline layout} описывает формат
\pause
\item \textit{Bind group} ссылается на конкретные данные (буфер/текстура/etc)
\pause
\item \textit{Render pipeline} описывает конкретный конвейер рендеринга (шейдеры + настройки примитивов + etc)
\end{itemize}
\end{frame}

\begin{frame}[fragile]
\frametitle{Bind groups + pipelines вместе}
\begin{itemize}
\item Разные bind группы могут использовать один и тот же bind group layout: например, layout материала объекта (набор текстур), у каждого объекта своя bind группа с конкретными текстурами
\pause
\item Разные render pipelines могут использовать один и тот же pipeline layout: например, ландшафт и 3D модели используют разные шейдеры, но один и тот же набор типов входных текстур и буферов (один и тот же \textit{layout})
\end{itemize}
\end{frame}

\begin{frame}[fragile]
\frametitle{Bind groups + pipelines: пример}
\usemintedstyle{lightbulb}
\begin{tcblisting}{listing engine=minted, minted language=cpp, minted options={fontsize=\scriptsize}, listing only, colback=codebg, colframe=codebg, rounded corners}
Bind group layout настроек камеры и освещения
  binding 0: uniform-буфер
  binding 1: storage-буфер

Bind group с настройками камеры и освещения
  binding 0: uniform-буфер с параметрами камеры
  binding 1: storage-буфер с источниками света

Bind group layout описания материала объекта
  binding 0: 2D текстура
  binding 1: 2D текстура
  binding 2: 2D текстура

Bind group с описанием материала объекта
  binding 0: текстура альбедо (цвет)
  binding 1: текстура нормалей
  binding 2: текстура материала (metallic-roughness)

Pipeline layout:
  group 0: настройки камеры и освещения
  group 1: описание материала объекта
\end{tcblisting}
\end{frame}

\begin{frame}[fragile]
\frametitle{Bind groups в шейдере}
\usemintedstyle{lightbulb}
\begin{tcblisting}{listing engine=minted, minted language=wgsl, minted options={fontsize=\scriptsize}, listing only, colback=codebg, colframe=codebg, rounded corners}
struct Camera {
  viewProjection: mat4x4f,
  position: vec3f,
  zRange: vec2f,
}

struct PointLight {
  position: vec3f,
  intensity: vec3f,
  falloff: vec3f,
}

@group(0) @binding(0) var<uniform> camera : Camera;
@group(0) @binding(1) var<storage, read> lights : array<PointLight>;

@group(1) @binding(0) var albedo_texture : texture_2d<f32>;
@group(1) @binding(1) var normal_texture : texture_2d<f32>;
@group(1) @binding(2) var material_texture : texture_2d<f32>;
@group(1) @binding(3) var linear_mipmap_sampler : sampler;
\end{tcblisting}
\end{frame}

\begin{frame}[fragile]
\frametitle{Bind groups + pipelines: рендеринг}
\usemintedstyle{lightbulb}
\begin{tcblisting}{listing engine=minted, minted language=rust, minted options={fontsize=\scriptsize}, listing only, colback=codebg, colframe=codebg, rounded corners}
render_pass.set_pipeline(my_pipeline)
render_pass.set_bind_group(0, camera_bind_group)

for object : objects {
  render_pass.set_bind_group(1, object.material_bind_group)
  render_pass.set_vertex_buffer(0, object.vertex_buffer)
  render_pass.set_index_buffer(object.index_buffer)
  render_pass.draw_indexed(object.index_count)
}
\end{tcblisting}
\end{frame}

\begin{frame}[fragile]
\frametitle{Uniform буферы}
\begin{itemize}
\item Буферы, данные которых можно передать как read-only в шейдеры
\pause
\item Uniform-данные гарантированно не могут меняться в рамках одной команды рисования (отсюда название \textit{uniform})
\pause
\item Для них у GPU есть специальные кеши памяти и отдельные (``скалярные'') инстукции для их чтения
\pause
\item \(\Longrightarrow\) Работают быстро :)
\end{itemize}
\end{frame}

\begin{frame}[fragile]
\frametitle{Uniform буферы}
\begin{itemize}
\item Такой буфер нужно создать с флагом \mintinline{cpp}|Uniform|
\pause
\item Тип в bind group layout: \mintinline{cpp}|buffer binding type = uniform|
\pause
\item Есть ограничения на размер буфера (точнее, его части), которую можно привязать к bind группе (обычно 64 Кб)
\end{itemize}
\end{frame}

\begin{frame}[fragile]
\frametitle{Uniform буферы vs immediates}
\begin{itemize}
\item У uniform-буфера размер гораздо больше, чем у immediates (в 512 раз!), зачем использовать immediates?
\pause
\item Uniform-буфер -- разновидность буфера, загрузить в него данные можно только через очередь \mintinline{cpp}|queue.write_buffer| (либо через асинхронный mapping, но там те же проблемы)
\pause
\item Последовательность команд на рисование отправляется в очередь целыми render pass'ами
\pause
\item \(\Longrightarrow\) Нельзя обновить данные uniform-буфера в середине render pass'а
\pause
\item \(\Longrightarrow\) Обычно в uniform буферах хранят данные, постоянные в течение одного кадра: параметры камеры, настройки глобального освещения
\end{itemize}
\end{frame}

\begin{frame}[fragile]
\frametitle{Uniform буферы: dynamic offset}
\begin{itemize}
\item Есть ограниченная поддержка динамического обновления uniform-буфера в виде \textit{dynamic offsets}
\pause
\item Можно включить dynamic offsets для конкретного буфера в bind group layout: \mintinline{cpp}|has_dynamic_offset = true|
\pause
\item Функция \mintinline{cpp}|render_pass.set_bind_group()| принимает массив \textit{dynamic offsets} -- по одному на каждый uniform-буфер в этой группе, для которых включен dynamic offset
\pause
\item Тогда будет читаться кусок буфера, начинающийся с этого offset'а
\pause
\item Сами offset'ы должны быть кратны значению \mintinline{cpp}|minUniformBufferOffsetAlignment| (по умолчанию -- 256)
\pause
\item \(\Longrightarrow\) Можно загрузить несколько наборов данных в uniform буфер и динамически переключаться между ними
\end{itemize}
\end{frame}

\begin{frame}[fragile]
\frametitle{Uniform буферы: шейдер}
\usemintedstyle{lightbulb}
\begin{tcblisting}{listing engine=minted, minted language=wgsl, minted options={fontsize=\scriptsize}, listing only, colback=codebg, colframe=codebg, rounded corners}
struct Camera {
  viewProjection: mat4x4f,
}

@group(0) @binding(0) var<uniform> camera : Camera;

struct Immediates {
  model: mat4x4f,
}

var<immediate> immediates;

@vertex
fn vertex_main(in: VertexIn) -> @builtin(position) vec4f {
  return camera.viewProjection
    * immediates.model
    * vec4f(in.position, 1.0);
}
\end{tcblisting}
\end{frame}

\begin{frame}[fragile]
\frametitle{Uniform буферы: инициализация}
\usemintedstyle{lightbulb}
\begin{tcblisting}{listing engine=minted, minted language=rust, minted options={fontsize=\scriptsize}, listing only, colback=codebg, colframe=codebg, rounded corners}
let bind_group_layout = device.create_bind_group_layout([
  BindGroupLayoutEntry {
    binding: 0,
    visibility: VERTEX,
    type: Buffer { type: Uniform },
  },
])

let uniform_buffer = device.create_buffer({
  size: 64,
  usage: CopyDst | Uniform,
})

let bind_group = device.create_bind_group(bind_group_layout, [
  BindGroupEntry {
    binding: 0,
    buffer: uniform_buffer,
  },
])
\end{tcblisting}
\end{frame}

\begin{frame}[fragile]
\frametitle{Uniform буферы: рендеринг}
\usemintedstyle{lightbulb}
\begin{tcblisting}{listing engine=minted, minted language=rust, minted options={fontsize=\scriptsize}, listing only, colback=codebg, colframe=codebg, rounded corners}
queue.write_buffer(uniform_buffer, &camera_matrix)

... 

render_pass.set_bind_group(0, bind_group)
render_pass.set_pipeline(...)
render_pass.set_immediates(...)
render_pass.draw(...)

...

queue.submit([command_buffer.finish()])
\end{tcblisting}
\end{frame}

\begin{frame}[fragile]
\frametitle{Текстуры}
\begin{itemize}
\item Текстуры -- изображения в памяти GPU (чаще -- несколько изображений)
\pause
\item Их можно читать из шейдеров
\pause
\item В них можно загружать пиксели из CPU
\pause
\item В них можно рисовать (вместо рисования в окно -- про это в следующий раз)
\pause
\item Поддерживают много специфичных для изображений операций
\end{itemize}
\end{frame}

\begin{frame}<1>[fragile,label=types]
\frametitle{Типы текстур}
\usemintedstyle{solarized-light}
\begin{itemize}
\item 1D: одномерная текстура (полоска пикселей)
\pause
\item 2D: двумерная текстура
\pause
\item 3D: трёхмерная текстура
\pause
\item 1D array: массив одномерных текстур одинакового размера (не поддерживается в WebGPU)
\pause
\item 2D array: массив двумерных текстур одинакового размера
\pause
\item Cubemap -- текстуры, хранящиеся как 6 граней куба
\pause
\item Cubemap array
\end{itemize}
\end{frame}

\begin{frame}[fragile]
\frametitle{Одномерная текстура}
\begin{itemize}
\item Градиенты цветов
\item Предподсчитанные значения функций (напр. для освещения)
\end{itemize}
\slideimage{texture_1d.png}
\end{frame}

\againframe<1-2>{types}

\begin{frame}[fragile]
\frametitle{Двумерная текстура}
\slideimage{texture_2d.jpg}
\end{frame}

\begin{frame}[fragile]
\frametitle{Двумерная текстура}
\begin{itemize}
\item Самый часто используемый тип
\item Любые числовые свойства поверхностей: цвет (альбедо), нормаль, тип материала
\item Промежуточные изображения при рендеринге: G-buffer, ID (visibility) buffer, HDR buffer
\item Карта высот ландшафта
\item Элементы UI
\item Предподсчитанные значения 2D функции
\item etc.
\end{itemize}
\end{frame}

\begin{frame}[fragile]
\frametitle{Карта нормалей}
\slideimage{normal_map.png}
\end{frame}

\begin{frame}[fragile]
\frametitle{G-buffer}
\slideimage{gbuffer.jpg}
\end{frame}

\begin{frame}[fragile]
\frametitle{Карта высот}
\slideimage{iceland_heightmap.png}
\end{frame}

\againframe<2-3>{types}

\begin{frame}[fragile]
\frametitle{Трёхмерная текстура}
\slideimage{texture_3d.png}
\end{frame}

\begin{frame}[fragile]
\frametitle{Трёхмерная текстура}
\begin{itemize}
\item Любые числовые свойства пространства: плотность, цвет, коэффициент рассеяния
\item Объёмные (volumetric) данные: дым, туман, и т.п.
\item Воксельные данные
\item Предподсчитанные значения 3D функции
\end{itemize}
\end{frame}

\againframe<3-4>{types}

\begin{frame}[fragile]
\frametitle{Массив одномерных текстур}
\begin{itemize}
\item Набор градиентов цветов
\end{itemize}
\slideimage{texture_1d_array.png}
\end{frame}

\againframe<4-5>{types}

\begin{frame}[fragile]
\frametitle{Массив двумерных текстур}
\begin{itemize}
\item Массив текстур для однотипных объектов, например, вокселей
\end{itemize}
\slideimage{texture_2d_array.png}
\end{frame}

\againframe<5-6>{types}

\begin{frame}[fragile]
\frametitle{Cubemap-текстура}
\begin{itemize}
\item Небо (skybox), отражения (environment map, reflection map), предпосчитанные значения функции на сфере
\end{itemize}
\slideimage{texture_cubemap.png}
\end{frame}

\againframe<6->{types}

\begin{frame}[fragile]
\frametitle{Текстуры: размер}
\begin{itemize}
\item Текстура -- массив (1D/2D/3D) \textit{пикселей} оперделённого \textit{размера} в определённом \textit{формате}
\pause
\item Размер -- целочисленный трёхмерный вектор, все компоненты больше нуля
\pause
\item Для array текстур последняя компонента размера -- количество элементов (\textit{слоёв, layers}) в массиве
\pause
\item Для cubemap текстур последняя компонента размера -- 6
\pause
\item Для cubemap array текстур последняя компонента размера -- 6 умножить на количество слоёв
\end{itemize}
\end{frame}

\begin{frame}[fragile]
\frametitle{Текстуры: формат пикселей}
\begin{itemize}
\item Форматы пикселей похожи на форматы вершин:
\begin{itemize}
\item \mintinline{cpp}|R8Unorm| -- одна компонента в 1 байт, беззнаковый нормированный (\([0\dots 255]\) превращается в \([0.0\dots 1.0]\))
\item \mintinline{cpp}|RG16Sint| -- две компоненты в 2 байта каждая, знаковый ненормированный (конверсия \mintinline{cpp}|int16| -> \mintinline{cpp}|float32|)
\item \mintinline{cpp}|RGBA32Float| -- четыре компоненты в 4 байта каждая, floating-point (значения попадают в шейдер как есть)
\item И т.п.
\end{itemize}
\pause
\item Есть особые форматы:
\begin{itemize}
\item \mintinline{cpp}|RGB10A2Unorm| -- три компоненты по 10 бит, одна в 2 бита, беззнаковые нормированные
\item \mintinline{cpp}|RG11B10Ufloat| -- две компоненты в 11-битном floating-point, одна в 10-битном
\item \mintinline{cpp}|RGB9E5Ufloat| -- три компоненты: отдельные мантиссы по 9 бит, общая экспонента в 5 бит
\end{itemize}
\end{itemize}
\end{frame}

\begin{frame}[fragile]
\frametitle{Текстуры: формат пикселей}
\begin{itemize}
\item Есть форматы для буфера глубины / stencil:
\begin{itemize}
\item \mintinline{cpp}|Stencil8| -- только 8-битный stencil буфер
\item \mintinline{cpp}|Depth16Unorm| -- 16-битный буфер глубины, беззнаковый нормированный
\item \mintinline{cpp}|Depth24Plus| -- 24-битный буфер глубины (может быть нормированным или floating-point)
\item \mintinline{cpp}|Depth24PlusStencil8| -- 24-битный буфер глубины + 8-битный stencil буфер
\item \mintinline{cpp}|Depth32Float| -- 32-битный floating-point буфер глубины
\item \mintinline{cpp}|Depth32FloatStencil8| -- 32-битный floating-point буфер глубины + 8-битный stencil буфер
\item И т.п.
\end{itemize}
\pause
\item Есть сжатые форматы: \mintinline{cpp}|BC*|, \mintinline{cpp}|ETC*|, \mintinline{cpp}|EACR*|, \mintinline{cpp}|ASTC*|
\pause
\item Есть sRGB-форматы: \mintinline{cpp}|RGBA8UnormSrgb|, и т.п.
\end{itemize}
\end{frame}

\begin{frame}[fragile]
\frametitle{Текстуры: создание}
\begin{itemize}
\item При создании текстуры нужно указать
\pause
\begin{itemize}
\item \mintinline{cpp}|dimension| -- 1D/2D/3D
\pause
\item \mintinline{cpp}|size| -- размер массива пикселей (по трём осям)
\pause
\item \mintinline{cpp}|usage| -- как будет использоваться текстура (набор флагов, аналогично buffer usage)
\pause
\item ...и ещё несколько вещей, о которых узнаем позже :)
\end{itemize}
\end{itemize}
\end{frame}

\begin{frame}[fragile]
\frametitle{Текстуры: usage}
\begin{itemize}
\item \mintinline{cpp}|CopySrc/CopyDst| -- будем копировать данные из текстуры / в текстуру
\pause
\item \mintinline{cpp}|TextureBinding| -- будем читать из шейдера как обычную текстуру (и добавлять в bind группы)
\pause
\item \mintinline{cpp}|StorageBinding| -- будем читать из шейдера как storage текстуру (и добавлять в bind группы)
\pause
\item \mintinline{cpp}|RenderAttachment| -- будем рисовать в текстуру с помощью растеризации
\pause
\item \mintinline{cpp}|TransientAttachment| -- render attachment, который не нужен после окончания render pass'а
\end{itemize}
\end{frame}

\begin{frame}[fragile]
\frametitle{Текстуры: views}
\begin{itemize}
\item Сама текстура -- только контейнер пикселей, чтобы её использовать в шейдере, нужно создать \textit{texture view}: \mintinline{cpp}|texture.create_view()|
\pause
\item Описывает, как \textit{интерпретировать} текстуру
\pause
\item Может читать 3D текстуру как 3D, 2D array, cubemap, или cubemap array
\pause
\item Может читать текстуру в другом формате (конвертировать из sRGB-форматов в обычные и наоборот)
\pause
\item Может читать только depth или stencil часть depth+stencil формата
\pause
\item Может читать конкретный диапазон слоёв array-текстуры или mipmap-уровней (о них позже)
\end{itemize}
\end{frame}

\begin{frame}[fragile]
\frametitle{Текстуры: bind group}
\begin{itemize}
\item В bind group layout'е:
\begin{itemize}
\item \mintinline{cpp}|view_dimension| -- 1D / 2D / 2D array / cubemap / etc
\item \mintinline{cpp}|sample_type| -- \mintinline{cpp}|Float| (для floating-point или нормированных текстур), \mintinline{cpp}|Sint| для знаковых целых, \mintinline{cpp}|Uint| для беззнаковых целых, \mintinline{cpp}|Depth| для глубины
\end{itemize}
\pause
\item В bind группе:
\begin{itemize}
\item \mintinline{cpp}|texture_view|
\end{itemize}
\end{itemize}
\end{frame}

\begin{frame}[fragile]
\frametitle{Текстуры: шейдер}
\begin{itemize}
\item В шейдере текстура представлена переменной с особым типом, соответствующим view dimension и sample type
\end{itemize}
\usemintedstyle{lightbulb}
\begin{tcblisting}{listing engine=minted, minted language=rust, minted options={fontsize=\scriptsize}, listing only, colback=codebg, colframe=codebg, rounded corners}
// T = f32 / i32 / u32
texture_1d<T>
texture_2d<T>
texture_2d_array<T>
texture_3d<T>
texture_cube<T>
texture_cube_array<T>
\end{tcblisting}
\end{frame}

\begin{frame}[fragile]
\frametitle{Тестуры: чтение}
\begin{itemize}
\item \mintinline{wgsl}|textureLoad(texture, coords, [array_layer], level)| -- возвращает \mintinline{glsl}|vec4f/vec4i/vec4u| в зависимости от типа текстуры (даже если в текстуре на самом деле меньше компонент)
\pause
\item \mintinline{wgsl}|coords| -- целочисленные координаты пикселя
\pause
\item \mintinline{wgsl}|array_layer| -- номер слоя для array-текстур
\pause
\item \mintinline{wgsl}|level| -- номер mipmap-уровня (о них позже)
\end{itemize}
\end{frame}

\begin{frame}[fragile]
\frametitle{Текстуры: чтение}
\usemintedstyle{solarized-light}
\begin{itemize}
\item Текстура хранит набор дискретных пикселей
\pause
\item Вычислять координаты пикселя -- неудобно, нужно знать размер текстуры
\pause
\item Хочется обращаться по неким (проинтерполированным из вершинного шейдера) floating-point координатам
\pause
\item Хочется сглаживать (интерполировать) значения между пикселями
\pause
\item \begin{math}\Longrightarrow\end{math} режимы фильтрации текстур
\end{itemize}
\end{frame}

\begin{frame}[fragile]
\frametitle{Текстуры: фильтрация}
\usemintedstyle{solarized-light}
\begin{itemize}
\item Режимы фильтрации:
\begin{itemize}
\item \mintinline{cpp}|Nearest| -- использовать ближайший к указанным координатам пиксель
\item \mintinline{cpp}|Linear| -- использовать линейную/билинейную/трилинейную интерполяцию между соседними 2/4/8 пикселями (для 1D/2D/3D соответственно)
\item \textbf{\alert{NB}}: Для array-текстур номер слоя \textit{всегда ближайший}!
\end{itemize}
\pause
\item Настраиваются отдельно для случая, когда
\begin{itemize}
\item Пиксель текстуры больше пикселя на экране (magnification) -- \mintinline{cpp}|mag_filter|
\item Пиксель текстуры меньше пикселя на экране (minification) -- \mintinline{cpp}|min_filter|
\end{itemize}
\end{itemize}
\end{frame}

\begin{frame}
\frametitle{Nearest vs Linear}
\slideimage{pikachu.png}
\end{frame}

\begin{frame}
\frametitle{Nearest vs Linear}
\slideimage{nearest_linear.png}
\end{frame}

\begin{frame}
\frametitle{Minification vs Magnification}
\slideimage{earth.png}
\end{frame}

\begin{frame}[fragile]
\frametitle{Mipmaps}
\begin{itemize}
\item Когда пиксель текстуры меньше пикселя на экране (minification), текстура выглядит плохо:
\begin{itemize}
\item \mintinline{cpp}|Nearest| -- часть пикселей текстуры не попадают на экран
\pause
\item \mintinline{cpp}|Linear| -- нормально до x2 уменьшения, потом пиксели тоже не попадают на экран
\end{itemize}
\pause
\item Решение: \textit{mipmap-уровни}
\end{itemize}
\end{frame}

\begin{frame}[fragile]
\frametitle{Mipmaps}
\begin{itemize}
\item Mipmaps -- уменьшенные копии текстуры для использования при \textit{минификации}
\pause
\item Для 2D текстуры \begin{math}W\times H\end{math} первый уровень имеет размер \begin{math}\left\lfloor\frac{W}{2}\right\rfloor\times \left\lfloor\frac{H}{2}\right\rfloor\end{math}, второй уровень имеет размер \begin{math}\left\lfloor\frac{W}{4}\right\rfloor\times \left\lfloor\frac{H}{4}\right\rfloor\end{math}, и т.д.
\pause
\item Последний mipmap-уровень имеет размер \begin{math}1\times 1\end{math}
\pause
\item Для 1D или 3D текстур -- аналогично по всем осям
\pause
\item Для array-текстур -- отдельный набор mipmap'ов для каждого слоя
\pause
\item Для cubemap-текстур -- отдельный набор mipmap'ов для каждой грани
\end{itemize}
\end{frame}

\begin{frame}
\frametitle{Mipmaps}
\slideimage{mipmaps.png}
\end{frame}

\begin{frame}[fragile]
\frametitle{Mipmaps}
\begin{itemize}
\item Количество mipmap-уровней указывается в поле \mintinline{cpp}|mip_level_count| при создании текстуры
\pause
\item Их пиксели нужно загрузить на GPU так же, как саму текстуру, либо сгенерировать уже на GPU (рендерингом или compute шейдерами)
\pause
\item Режимы фильтрации mipmap-уровней -- \mintinline{cpp}|mipmap_filter|:
\begin{itemize}
\item \mintinline{cpp}|Nearest| -- выбирается ближайший mipmap-уровень
\pause
\item \mintinline{cpp}|Linear| -- берутся пиксели с двух ближайших уровней и линейно интерполируются
\end{itemize}
\end{itemize}
\end{frame}

\begin{frame}[fragile]
\frametitle{Би/трилинейная фильтрация}
\begin{itemize}
\item Размерность текстуры (1D/2D/3D) и тип фильтрации влияют на общее число линейных интерполяций, отсюда названия \textit{билинейная} фильтрация и \textit{трилинейная} фильтрация
\pause
\item Часто под трилинейной фильтрацией имеют в виду конкретно 2D текстуру с \mintinline{cpp}|min_filter = Linear| и \mintinline{cpp}|mipmap_filter = Linear|
\end{itemize}
\end{frame}

\begin{frame}[fragile]
\frametitle{Анизотропная фильтрация}
\begin{itemize}
\item Текстура после проецирования на экран может быть нормального размера по одной оси, и маленькая по другой
\pause
\item Большой mipmap уровень хорош для одной оси, маленький -- для другой; ни один уровен не будет выглядеть хорошо
\pause
\item Решение: \textit{анизотропная фильтрация} -- читать текстуру в нескольких местах сразу и усреднять (поддерживается аппаратно)
\pause
\item Включается настройкой \mintinline{cpp}|max_anisotropy| -- от 1 до 16 (количество чтений из текстуры, которые делает GPU)
\end{itemize}
\end{frame}

\begin{frame}
\frametitle{Анизотропная фильтрация}
\slideimage{anisotropy.png}
\end{frame}

\begin{frame}[fragile]
\frametitle{Wrapping mode}
\begin{itemize}
\item При чтении из текстуры координаты могут выходить за пределы диапазона
\pause
\item При \mintinline{cpp}|Linear|-фильтрации для точек у границы текстуры может не найтись достаточное количество соседних пикселей
\pause
\item \begin{math}\Longrightarrow\end{math} \mintinline{cpp}|address mode|:
\begin{itemize}
\item \mintinline{cpp}|ClampToEdge| -- брать пиксель с границы
\item \mintinline{cpp}|ClampToBorder| -- брать пиксель фиксированного цвета
\item \mintinline{cpp}|Repeat| -- повторять текстуру
\item \mintinline{cpp}|MirrorRepeat| -- повторять текстуру, отражая её на каждый повтор
\end{itemize}
\pause
\item Настраивается отдельно для каждой текстурной координаты (U, V, W)
\end{itemize}
\end{frame}

\begin{frame}
\frametitle{Wrapping}
\slideimage{wrapping.png}
\end{frame}
\begin{frame}
\frametitle{Wrapping}
\slideimage{wrapping2.jpg}
\end{frame}

\begin{frame}[fragile]
\frametitle{Samplers}
\begin{itemize}
\item Все настройки фильтрации и wrapping'а хранятся не в самой текстуре, а в отдельном объекте -- \textit{сэмплере}
\pause
\item \mintinline{cpp}|device.create_sampler({ addressModeU = Repeat, ... })|
\pause
\item Его тоже нужно добавлять в bind группы (и, соответственно, в layout'ы)
\pause
\item \textbf{\alert{NB}}: В некоторых API (OpenGL, Vulkan) можно объеденить текстуру и сэмплер в один объект для удобства
\end{itemize}
\end{frame}

\begin{frame}[fragile]
\frametitle{Текстуры: чтение}
\begin{itemize}
\item Процесс чтения текстуры с использованием сэмплера (и его настроек) называется \textit{сэмплингом} текстуры
\pause
\item Для сэмплинга есть отдельная функция: \mintinline{wgsl}|textureSample(texture, sampler, coords, array_layer)|
\pause
\item Здесь, \mintinline{wgsl}|coords| -- не номер пикселя, а \textit{текстурные координаты}:
\pause
\begin{itemize}
\item Для 1D/2D/3D текстур \mintinline{wgsl}|float/vec2f/vec3f| -- \textit{нормированные} координаты, [0..1] по каждой оси
\pause
\item Для cubemap текстур \mintinline{wgsl}|vec3f| -- вектор направления из центра куба: возвращается значение из пересечения луча с поверхностью куба
\pause
\item Частое название в случае 2D текстур: \textit{UV-координаты}
\end{itemize}
\end{itemize}
\end{frame}

\begin{frame}[fragile]
\frametitle{Текстурные координаты}
\begin{itemize}
\item \mintinline{wgsl}|texture_2d|: \mintinline{wgsl}|vec2(0.6, 0.5)| -- пиксель чуть правее центра текстуры
\pause
\item \mintinline{wgsl}|texture_cube|: \mintinline{wgsl}|vec3(0.1, 1.0, -0.1)| -- пиксель на +Y грани куба, около центра грани
\end{itemize}
\end{frame}

\begin{frame}[fragile]
\frametitle{Выбор mipmap уровня: shader derivatives}
\begin{itemize}
\item Как GPU понимает, какой выбрать mipmap-уровень, и происходит ли minification/magnification? \pause По значениям в \textit{соседних пикселях}
\pause
\item Фрагментный шейдер всегда запускается на группах 2x2 пикселя (даже если часть из них не были растеризованы)
\pause
\item В GLSL есть функции \mintinline{wgsl}|dpdx/dpdy| -- вычисляют разницу некой величины для пары соседних пикселей в группе 2x2
\pause
\begin{itemize}
\item Например, \mintinline{wgsl}|dpdx(value)| -- разница между значением \mintinline{wgsl}|value| в пикселе справа и значением в пикселе слева
\end{itemize}
\pause
\item По \mintinline{wgsl}|dpdx(coords)| и \mintinline{wgsl}|dpdy(coords)| вычисляется mipmap-уровень (где \mintinline{wgsl}|coords| -- текстурные координаты)
\pause
\item Интерпретируя \mintinline{wgsl}|dpdx/dpdy| как аппроксимации производных, можно делать много интересных трюков: восстанавливать нормаль к поверхности по z-буферу, рисовать линии фиксированной толщины даже с перспективной проекцией, и т.п.
\end{itemize}
\end{frame}

\begin{frame}[fragile]
\frametitle{Текстуры: загрузка данных}
\begin{itemize}
\item \mintinline{cpp}|queue.write_texture(copy_info, data_layout, data, size)|
\pause
\item \mintinline{cpp}|data_layout| -- описывает, как организованы байты текстуры в памяти:
\pause
\begin{itemize}
\item \mintinline{cpp}|bytes_per_row| -- количество байт на одну строку текстуры (может включать padding в конце строки)
\pause
\item \mintinline{cpp}|rows_per_image| -- количество строк на один 2D слой (имеет смысл для 3D текстур)
\end{itemize}
\pause
\item \mintinline{cpp}|copy_info| -- описывает, куда нужно записать пиксели:
\pause
\begin{itemize}
\item \mintinline{cpp}|texture| -- текстура, куда загружаем пиксели
\pause
\item \mintinline{cpp}|mip_level| -- в какой mipmap-уровень
\pause
\item \mintinline{cpp}|origin| -- координаты левого нижнего угла (можно загрузить часть текстуры в произвольном месте)
\end{itemize}
\end{itemize}
\end{frame}

\begin{frame}[fragile]
\frametitle{Текстуры: загрузка данных}
\begin{itemize}
\item \mintinline{cpp}|queue.copy_buffer_to_texture(...)| -- скоприровать данные из буфера в текстуру
\pause
\item Полезно для асинхронной загрузки текстур (через mapped buffer), для сжатия текстур на GPU
\pause
\item Для этой функции \mintinline{cpp}|bytes_per_row| должен быть кратен 256
\end{itemize}
\end{frame}

\begin{frame}[fragile]
\frametitle{Как использовать текстуру для 3D-модели?}
\begin{itemize}
\item Обычно каждая вершина 3D-модели хранит нормированные 2D текстурные координаты (\textit{UV-развёртка})
\pause
\item Они передаются из вершинного шейдера во фрагментный и интерполируются
\pause
\item Интерполированные значения текстурных координат используются во фрагментном шейдере для сэмплинга текстуры
\end{itemize}
\end{frame}

\begin{frame}[fragile]
\frametitle{Как использовать текстуру для 3D-модели?}
\slideimage{uvs.png}
\end{frame}

\begin{frame}[fragile]
\frametitle{Форматы изображений}
\begin{itemize}
\item \mintinline{cpp}|queue.write_texture()| ожидает на вход массив готовых пикселей
\pause
\item Обычно изображения хранятся в более сложных форматах: PNG, JPG, и т.д.
\pause
\item Есть формат Netpbm, хранящий сырые пиксели и простой для чтения/записи
\pause
\item Есть библиотеки для чтения разных форматов -- libpng, libjpeg
\pause
\item \mintinline{cpp}|stb_image.h| -- single-header библиотека для чтения разных форматов
\pause
\item Boost GIL (Generic Image Library)
\end{itemize}
\end{frame}

\begin{frame}[fragile]
\frametitle{Текстура: полный код (инициализация)}
\usemintedstyle{lightbulb}
\begin{tcblisting}{listing engine=minted, minted language=rust, minted options={fontsize=\scriptsize}, listing only, colback=codebg, colframe=codebg, rounded corners}
texture = device.create_texture({
  dimension: 2D,
  size: [1024, 1024, 1],
  usage: CopyDst | RenderAttachment,
})

queue.write_texture(...)

sampler = device.create_sampler({
  addressModeU: Repeat,
  minFilter: Linear,
  ...
})
\end{tcblisting}
\end{frame}

\begin{frame}[fragile]
\frametitle{Текстура: полный код (инициализация)}
\usemintedstyle{lightbulb}
\begin{tcblisting}{listing engine=minted, minted language=rust, minted options={fontsize=\scriptsize}, listing only, colback=codebg, colframe=codebg, rounded corners}
bind_group_layout = device.create_bind_group_layout([
  {
    binding: 0,
    texture: { Float, 2D },
  },
  {
    binding: 1,
    sampler: { Filtering },
  },
])

bind_group = device.create_bind_group(bind_group_layout, [
  {
    binding: 0,
    texture_view: texture.create_view(),
  },
  {
    binding: 1,
    sampler: sampler,
  },
])
\end{tcblisting}
\end{frame}

\begin{frame}[fragile]
\frametitle{Текстура: полный код (рендеринг)}
\usemintedstyle{lightbulb}
\begin{tcblisting}{listing engine=minted, minted language=rust, minted options={fontsize=\scriptsize}, listing only, colback=codebg, colframe=codebg, rounded corners}
render_pass.set_bind_group(0, bind_group)
render_pass.set_pipeline(my_pipeline)
render_pass.set_vertex_buffer(...)
render_pass.set_index_buffer(...)
render_pass.draw_indexed(...)
\end{tcblisting}
\end{frame}

\begin{frame}[fragile]
\frametitle{Текстура: полный код (шейдер)}
\usemintedstyle{lightbulb}
\begin{tcblisting}{listing engine=minted, minted language=wgsl, minted options={fontsize=\scriptsize}, listing only, colback=codebg, colframe=codebg, rounded corners}
@group(0) @binding(0) var color_texture: texture_2d<f32>;
@group(0) @binding(1) var linear_sampler: sampler;

struct VertexIn {
  @location(0) position: vec3f,
  @location(1) texcoord: vec2f,
}

struct VertexOut {
  @builtin(position) clip_position: vec4f,
  @location(0) texcoord: vec2f,
}

@vertex
fn vertex_main(in: VertexIn) -> VertexOut {
  return VertexOut(matrix * vec4f(in.position, 1.0), in.texcoord);
}

@fragment
fn fragment_main(in: VertexOut) -> @location(0) vec4f {
  return textureSample(color_texture, linear_sampler, in.texcoord);
}
\end{tcblisting}
\end{frame}

\begin{frame}[fragile]
\frametitle{Ссылки}
\begin{itemize}
\item \href{https://www.khronos.org/opengl/wiki/Texture}{\nolinkurl{khronos.org/opengl/wiki/Texture}}
\item \href{https://www.khronos.org/opengl/wiki/Sampler_Object}{\nolinkurl{khronos.org/opengl/wiki/Sampler\_Object}}
\item \href{https://www.khronos.org/opengl/wiki/Image_Format}{\nolinkurl{khronos.org/opengl/wiki/Image\_Format}}
\item \href{https://learnopengl.com/Getting-started/Textures}{\nolinkurl{learnopengl.com/Getting-started/Textures}}
\item \href{http://www.opengl-tutorial.org/beginners-tutorials/tutorial-5-a-textured-cube}{\nolinkurl{opengl-tutorial.org/beginners-tutorials/tutorial-5-a-textured-cube}}
\item \href{https://open.gl/textures}{\nolinkurl{open.gl/textures}}
\end{itemize}
\end{frame}

\end{document}